\documentclass[UTF8,12pt]{ctexart}
\usepackage[paperwidth=240mm,paperheight=200mm,margin=14mm,top=8mm]{geometry}
\usepackage{amsmath,amssymb,mathtools}
\usepackage{xcolor}
\pagestyle{empty}
\setlength{\parindent}{0pt}
\setlength{\parskip}{0pt}
\newcommand{\qn}[1]{\textbf{\large #1}\ }
\xeCJKDeclareCharClass{CJK}{"2460 -> "24FF}
\begin{document}
\qn{1999.1}
曲线 $\begin{cases} x = e^t \sin 2t \\ y = e^t \cos t \end{cases}$ 在点 $(0,1)$ 处的法线方程为 $\_\_\_\_\_\_$.

\vfill
\newpage
\qn{1999.2}
设函数 $y = y(x)$ 由方程 $\ln(x^2 + y) = x^3 y + \sin x$ 确定，则 $\left.\dfrac{dy}{dx}\right|_{x=0} = \_\_\_\_\_\_.$

\vfill
\newpage
\qn{1999.3}
$\int \dfrac{x + 5}{x^2 - 6x + 13} \, dx = \_\_\_\_\_\_.$

\vfill
\newpage
\qn{1999.4}
函数 $y = \dfrac{x^2}{\sqrt{1 - x^2}}$ 在区间 $\left[\dfrac{1}{2}, \dfrac{\sqrt{3}}{2}\right]$ 上的平均值为 $\_\_\_\_\_\_.$

\vfill
\newpage
\qn{1999.5}
微分方程 $y'' - 4y = e^{2x}$ 的通解为 $\_\_\_\_\_\_.$

\vfill
\newpage
\qn{1999.6}
设 $f(x) = \begin{cases} \dfrac{1 - \cos x}{\sqrt{x}}, & x > 0 \\ x^2 g(x), & x \leqslant 0 \end{cases}$，其中 $g(x)$ 是有界函数，则 $f(x)$ 在 $x = 0$ 处（　）
\\
\noindent (A) 极限不存在.　　　　　(B) 极限存在，但不连续.
\\
\noindent (C) 连续，但不可导.　　　(D) 可导.

\vfill
\newpage
\qn{1999.7}
设 $\alpha(x) = \int_0^{5x} \dfrac{\sin t}{t} \, dt$，$\beta(x) = \int_0^{\sin x} (1 + t)^{\frac{1}{t}} \, dt$，则当 $x \to 0$ 时，$\alpha(x)$ 是 $\beta(x)$ 的（　）
\\
\noindent (A) 高阶无穷小.　　　　　(B) 低阶无穷小.
\\
\noindent (C) 同阶但不等价的无穷小.　　　　　(D) 等价无穷小.

\vfill
\newpage
\qn{1999.8}
设 $f(x)$ 是连续函数，$F(x)$ 是 $f(x)$ 的原函数，则（　）
\\
\noindent (A) 当 $f(x)$ 是奇函数时，$F(x)$ 必是偶函数.
\\
\noindent (B) 当 $f(x)$ 是偶函数时，$F(x)$ 必是奇函数.
\\
\noindent (C) 当 $f(x)$ 是周期函数时，$F(x)$ 必是周期函数.
\\
\noindent (D) 当 $f(x)$ 是单调增函数时，$F(x)$ 必是单调增函数.

\vfill
\newpage
\qn{1999.9}
“对任意给定的 $\varepsilon \in (0,1)$，总存在正整数 $N$，当 $n \geqslant N$ 时，恒有 $|x_n - a| \leqslant 2\varepsilon$” 是数列 $\{x_n\}$ 收敛于 $a$ 的（　）
\\
\noindent (A) 充分条件但非必要条件.　　　　　(B) 必要条件但非充分条件.
\\
\noindent (C) 充分必要条件.　　　　　(D) 既非充分条件又非必要条件.

\vfill
\newpage
\qn{1999.10}
记行列式 $\begin{vmatrix} x - 2 & x - 1 & x - 2 & x - 3 \\ 2x - 2 & 2x - 1 & 2x - 2 & 2x - 3 \\ 3x - 3 & 3x - 2 & 4x - 5 & 3x - 5 \\ 4x & 4x - 3 & 5x - 7 & 4x - 3 \end{vmatrix}$ 为 $f(x)$，则方程 $f(x) = 0$ 的根的个数为（　）
\\
\noindent (A) 1.　　　　　(B) 2.　　　　　(C) 3.　　　　　(D) 4.
求 $\lim_{x \to 0} \dfrac{\sqrt{1 + \tan x} - \sqrt{1 + \sin x}}{x \ln(1 + x) - x^2}$.
计算 $\int_1^{+\infty} \dfrac{\arctan x}{x^2} \, dx$.
求初值问题 $\begin{cases} (y + \sqrt{x^2 + y^2}) \, dx - x \, dy = 0 \quad (x > 0) \\ \left.y\right|_{x=1} = 0 \end{cases}$ 的解.
为清除井底的污泥，用缆绳将抓斗放入井底，抓起污泥后提出井口，已知井深 $30\,\text{m}$，抓斗自重 $400\,\text{N}$，缆绳每米重 $50\,\text{N}$，抓斗抓起的污泥重 $2000\,\text{N}$，提升速度为 $3\,\text{m/s}$. 在提升过程中，污泥以 $20\,\text{N/s}$ 的速率从抓斗缝隙中漏掉. 现将抓起污泥的抓斗提升至井口，问克服重力需作多少焦耳的功？
（说明：① $1\,\text{N} \times 1\,\text{m} = 1\,\text{J}$；$\text{m},\text{N},\text{s},\text{J}$ 分别表示米、牛顿、秒、焦耳. ② 抓斗的高度及位于井口上方的缆绳长度忽略不计.）
已知函数 $y = \dfrac{x^3}{(x - 1)^2}$，求

\vfill
\newpage
\qn{1999.11}
函数的增减区间及极值；

\vfill
\newpage
\qn{1999.12}
函数图形的凹凸区间及拐点；

\vfill
\newpage
\qn{1999.13}
函数图形的渐近线.
设函数 $f(x)$ 在闭区间 $[-1,1]$ 上具有三阶连续导数，且 $f(-1) = 0$，$f(1) = 1$，$f'(0) = 0$，证明：在开区间 $(-1,1)$ 内至少存在一点 $\xi$，使 $f'''(\xi) = 3$.
设函数 $y(x) \, (x \geqslant 0)$ 二阶可导，且 $y'(x) > 0$，$y(0) = 1$. 过曲线 $y = y(x)$ 上任意一点 $P(x,y)$ 作该曲线的切线及 $x$ 轴的垂线，上述两直线与 $x$ 轴所围成的三角形的面积记为 $S_1$，区间 $[0,x]$ 上以 $y = y(x)$ 为曲边的曲边梯形面积记为 $S_2$，并设 $2S_1 - S_2$ 恒为 $1$，求此曲线 $y = y(x)$ 的方程.
设 $f(x)$ 是区间 $[0, +\infty)$ 上单调减少且非负的连续函数，$a_n = \sum_{k=1}^n f(k) - \int_1^n f(x) \, dx \quad (n = 1,2,\cdots)$，证明数列 $\{a_n\}$ 的极限存在.
设矩阵 $A = \begin{pmatrix} 1 & 1 & -1 \\ -1 & 1 & 1 \\ 1 & -1 & 1 \end{pmatrix}$，矩阵 $X$ 满足 $A^* X = A^{-1} + 2X$，其中 $A^*$ 是 $A$ 的伴随矩阵，求矩阵 $X$.
设向量组 $\alpha_1 = (1,1,1,3)^T$，$\alpha_2 = (-1,-3,5,1)^T$，$\alpha_3 = (3,2,-1,p+2)^T$，$\alpha_4 = (-2,-6,10,p)^T$.

\vfill
\newpage
\qn{1999.14}
$p$ 为何值时，该向量组线性无关？并在此时将向量 $\alpha = (4,1,6,10)^T$ 用 $\alpha_1,\alpha_2,\alpha_3,\alpha_4$ 线性表示；

\vfill
\newpage
\qn{1999.15}
$p$ 为何值时，该向量组线性相关？并在此时求出它的秩和一个极大线性无关组.

\vfill
\newpage
\end{document}